% !TeX root = main.tex
\section{模拟与编程}

\subsection{数字孪生}

数字孪生（Digitaler Zwilling）的基本设想是为物理世界中的每一项\textbf{资产}都建立一份与之相对应、并在其整个生命周期内持续维护的数字化表征，使得原本只能在实物上进行的操作能够转移到数字侧。

数字孪生并不以完整复刻物理实体为目标。按照定义，数字孪生是\textbf{一种详尽程度足以满足一组既定应用场景之需求的数字化表征}，其边界并非由物理实体本身划定，而是由使用目的反向确定。这意味着同一项资产在不同用途下可以对应不同的孪生。

与之相对，数字影子（Digitaler Schatten）指\textbf{足够精确地反映生产、开发及相邻领域中各类过程的映像}，其目的在于为一切相关数据建立一个\textbf{具备实时处理能力的评估基础}，具体包括对必要数据格式、数据选择及数据粒度层级的描述。

数字影子侧重于由物理侧向数字侧的采集与呈现，构成的是一个\textbf{观测与评估}的基础。

数字孪生要求与其物理对应物之间保持\textbf{双向的数据关联}，数字侧得出的结论能够反作用于物理侧的运行与决策。

若数字孪生的发展延续既有轨迹，未来它很可能\textbf{与管理壳（Verwaltungsschale）这一概念趋同}，而在工业4.0的语境中，管理壳这一术语正逐渐成为优先使用的说法。

\subsection{仿真}

要准确理解仿真，需要先厘清一组彼此关联的核心概念。首先是系统本身：一个系统是由若干相互关联的组件构成的、有边界的排列，其构成要素包括系统边界、系统的输入与输出量、组件与子系统及系统元素、运行结构、组成结构、状态量与状态转换；若一个系统的组成部分本身又是相对独立的系统，则称其为“系统之系统”。以某个可重构装配单元为例，其典型的输入包括零件、订单数据与能量，输出则包括产品、协议记录与排放物，内部则由操作员、控制台、工件、相机、单元控制、六足平台、环境、机器人与夹持器等要素及其各自的状态变量共同构成。

\begin{figure}[htbp]
  \centering
  \includegraphics[width=0.7\textwidth]{simulationcycle.png}
  \caption{“现实系统—概念模型—仿真模型—可实验模型”四方框循环图}
  \label{fig:simulationcycle}
\end{figure}

在系统之上，模型是对一个规划中或已存在系统及其过程的简化重现，重现所依托的是另一种概念性或实体性的系统。模型与原型之间的差异，仅限于研究目标所划定的容差范围之内，且这种差异只体现在与研究相关的特性上。模型可以在不同的抽象层次上（分析、设计、实现）、针对不同的系统组成部分或专业领域（工艺技术、机械、电气、自动化、信息技术）、围绕不同的关注点（安全性、性能、承载能力）、面向不同的任务（测试、投入使用）来表征系统，其类型既可以是概括性的理论，也可以是针对具体对象或系统的专门描述，模型所应具备的关键属性则包括映射性、抽象性与实用性。

在从现实系统到可执行仿真的路径中，存在几个依次递进的模型形态。概念模型是对系统的抽象、简化表征，应当采用使用者的术语来表达，其本身并不可执行，而是需要经过形式化处理，转化为仿真模型；概念模型描述了仿真模型的目标、输入输出、内容、假设与简化，是一种独立于具体软件的描述。仿真模型则是对一个规划中或已存在系统的正式的、面向仿真的具体重现，其构建方式多种多样，既可以直接使用编程语言（灵活但容易出错），也可以使用图形化或文本化的专用仿真语言（更为专门化，能够降低工作量与出错概率），还可以借助仿真工具自身的用户界面来完成，无论采取哪种方式，构建过程都需要关于目标系统足够充分且有说服力的数据支撑。仿真本身，则是在实验模型中重现系统及其动态过程，以获得可迁移到现实中的认识，其中系统过程会随时间演化；从更宽泛的意义上说，仿真也可以理解为围绕一个仿真模型开展有针对性实验的准备、执行与评估这一整套活动。承载这一切的软件工具被称为仿真工具或仿真器，它使得对系统动态行为及其过程加以重现和执行的模型得以被创建并运行起来。

上述各类模型之间的界限并非泾渭分明，而是逐渐过渡的：以机电一体化系统的模块化分解为例，机械、控制与电子这几个专业领域各自都会经历从拓扑模型描述、物理模型描述、数学模型描述到行为模型描述的过程，其中拓扑与物理层面的描述大体对应概念模型，数学层面的描述对应仿真模型，而行为层面的描述则对应可实验的模型，最终通过离散时间的连接与调度机制耦合在一起。将这一系列环节串联起来看，仿真本身是一个完整的过程：从明确研究的问题、目标与需求出发，经过概念建模（抽象）、形式化（转化为特定仿真语言下的仿真模型）、实现（转化为可执行代码构成的可实验模型），到实际开展仿真实验、获得仿真数据，最终对结果加以解读与整理；在这一过程的各个环节之间，还需要通过验证确认模型与其上一阶段的表征是否一致，并通过确认（validation）检验最终的仿真结果是否能够回答最初提出的问题。

在从模型走向可用结论的过程中，还需要区分两个容易混淆的概念：\textbf{验证（Verifikation）}是检查系统或模型是否满足事先\textbf{规定的技术要求}，也就是"有没有把东西做对"；\textbf{确认（Validierung）}则是检查系统或模型是否真正\textbf{达成了预期的使用目的}，也就是"有没有做对的东西"。就仿真模型而言，验证关心的是模型是否被正确地实现——方程、参数与程序有无差错；确认关心的则是模型是否足够贴近现实、能否回答最初提出的那个问题。

就今天的生产技术而言，仿真技术的典型应用可以按所考察的对象分为四类，每一类各有其相适配的方法（图 \ref{fig:simulationcatagory}）：
\begin{itemize}
  \item \textbf{物流系统}：关注物料流与订单流的时序、排队与产能瓶颈，通常采用\textbf{事件离散仿真}；
  \item \textbf{生产系统}：关注设备布局、可达性与工位之间的配合，通常采用\textbf{三维仿真}与三维可视化；
  \item \textbf{复杂机电系统}：关注多个刚体在关节耦合下的动态行为，通常采用\textbf{多体动力学}；
  \item \textbf{物理参数}：关注连续体内部的变形、应力与温度分布，通常采用\textbf{有限元方法}。
\end{itemize}
四者的差别本质上在于抽象层次的高低：从只关心事件发生的先后顺序，到关心几何与运动，再到关心力与力矩，直至关心材料内部的连续场分布。

\begin{figure}[htbp]
  \centering
  \includegraphics[width=0.75\textwidth]{simulationcatagory.png}
  \caption{生产技术中仿真的典型应用领域及其适配方法}
  \label{fig:simulationcatagory}
\end{figure}

\paragraph{面向机器人领域的专业仿真方法}

在机器人及相关机电系统的仿真中，\textbf{运动学（Kinematik）}与\textbf{动力学（Dynamik）}是两类基础而又彼此互补的方法。运动学描述的是刚体的运动可能性与实际运动，借助适当的坐标——例如位置、速度与加速度作为时间的函数——来刻画点与刚体的几何及解析运动状态，\textbf{它并不考虑引发运动的力与力矩}，而只关注位置、速度、加速度本身，典型应用包括描述机器人臂的位姿，或推导驱动机械臂所需的运动方程。刚体动力学则更进一步，研究质点、质点系、刚体与刚体系统在\textbf{受到力与力矩作用下}、并遵循运动学规律约束的运动情况；作为动力学的一个特殊而重要的情形，刚体动力学假设物体本身不可变形，用于描述自由运动或彼此耦合的多个刚体的动态行为，其理论基础是牛顿-欧拉与拉格朗日方法中关于能量和动量守恒的定律，刚体之间的耦合关系或接触所产生的约束，最终归结为整个系统的运动方程。与刚体动力学形成对照的是有限元方法：它是求解偏微分方程的一种数值方法，在时间之外还对空间进行离散，即将连续体划分为有限个小单元构成的网格，进而得到一个通常规模庞大（其规模取决于单元数目，对于线性偏微分方程则为线性）的方程组，并通过直接或迭代方式求解；有限元方法常用于描述固体中的静态与动态问题，例如变形、应力或温度分布，通过将部件置于预期的受载状况之下，识别出局部的薄弱环节，其基本原理可以概括为系统刚度矩阵（包含材料参数，例如弹性模量）联系外部作用与节点位移这一关系式。

不同的仿真工具可以依据"单一领域仿真—多领域仿真"以及"跨领域建模—仿真方法耦合"这两个维度加以归类：有限元类工具与数值数学类工具偏向单领域仿真，键合图（Bondgraph）类工具偏向跨领域建模，而多体仿真类工具则更侧重仿真方法之间的耦合，处于二者之间的还有各类耦合接口。键合图是一种跨领域描述物理系统的通用手段，其中节点代表物理元素或子系统，边则代表二者之间的相互作用（即所谓的"键"），其背后统一的原理在于：力与速度之积、电压与电流之积、温度与熵流之积、压强与体积流之积，均等于功率，从而把不同物理领域中形式各异的现象，纳入同一套描述框架之下。除此之外，围绕铸造、锻造、铣削、钻孔、焊接、钎焊等各类具体制造工艺，还存在大量面向具体应用的专门仿真方法，其种类远不止于此处所述。

\subsection{最大坐标系下的刚体动力学}

在刚体动力学的具体建模中，存在两种典型的坐标选择方式。

其一是最大坐标方法：动力学系统中\textbf{每个刚体各自拥有六个自由度}，\textbf{这些自由度可以通过约束条件加以限制}，\textbf{每一个约束条件都会对应产生一个约束力}，这一思路是牛顿-欧拉方法的典型做法——以一台六轴工业机器人为例，在最大坐标方法下，其对应的方程组维度为三十。

其二是最小坐标方法：\textbf{运动方程只针对系统实际的自由度建立}，\textbf{无需引入约束条件}，因而也\textbf{无需计算约束力}，这是拉格朗日方法的典型做法——同样以六轴工业机器人为例，在最小坐标方法下，其对应的方程组维度仅为六。

尽管最小坐标方法在方程规模上更为简洁，但从多个实际场景出发，仍然能够看出采用最大坐标方法的必要性：当约束条件（例如关节或碰撞接触）的数量随时间变化、当系统中存在并联运动结构、当系统具有自由运动的基座（例如卫星）、又或者需要求解动态支反力时，针对每一类具体情形单独设计求解方案往往并不经济，而最大坐标方法提供了一种统一的、可以涵盖各类情形的处理途径。

刚体动力学最大坐标方法的出发点，是牛顿-欧拉关于动量与角动量的定理：动量随时间的变化率等于所受合外力，角动量随时间的变化率等于所受合外力矩，后者具体表示为惯量张量与角加速度之积再加上角速度与惯量张量及角速度叉乘项之和。由此可以反解出加速度与角加速度的表达式，分别正比于外力除以质量，以及惯量张量的逆作用于外力矩与陀螺项之差。若需要同时仿真 $n$ 个自由运动的刚体，则可以把上述方程整理为矩阵形式：各刚体的平动与转动加速度所构成的向量，等于一个由各刚体质量矩阵与惯量张量之逆构成的分块对角矩阵，乘以各刚体所受外力与外力矩（含陀螺项）构成的向量。

为了在仿真中推进系统状态，需要对上述加速度-力关系做两次数值积分。以简单的欧拉积分法为例，加速度可以用速度的差商近似，从而得到下一时刻的速度等于当前速度加上积分步长乘以质量逆矩阵与外力之积；类似地，位置可以用速度的差商表示，从而得到下一时刻的位置等于当前位置加上积分步长乘以当前速度；这两步积分同样可以整理为统一的矩阵形式。

在现实的动力学系统中，刚体通常无法自由运动，而是通过关节相互连接，从而使得各刚体相对彼此的运动受到限制，这类限制在刚体动力学仿真中称为约束条件，其一般形式可以写成一个关于位置和时间的函数等于零向量。对这一位置层面的约束条件求导，便得到速度层面的约束条件，其中出现的雅可比矩阵 $J$ 正是约束条件相对于状态变量的偏导数。

根据达朗贝尔原理，对于满足约束条件、发生在无穷小且不耗费时间的虚位移而言，雅可比矩阵与速度向量之积必须为零。约束条件所引发的约束力，其作用是保证刚体只能沿着允许的轨迹运动；依据达朗贝尔原理，\textbf{约束力不做虚功}，即约束力与速度向量之积恒为零，若将约束力选取为与约束条件正交（即约束力等于雅可比矩阵的转置乘以待定的拉格朗日乘子），便自动满足达朗贝尔原理，这一点也可以通过代入验证：约束力与速度向量的内积可以改写为拉格朗日乘子与雅可比矩阵和速度向量之积的乘积，而后者恰好为零。以一个平面上受约束运动的质点为例，其约束条件为 $y=2x$，对应的雅可比矩阵为 $(2, -1)$，可以验证约束力方向与该约束边界正交，且与外力方向相容的速度分量满足该约束关系。

将约束力纳入运动方程的积分之中，便得到最终用于仿真一个刚体系统的方程组：下一时刻的加速度状态等于当前状态加上积分步长乘以质量逆矩阵与外力及约束力之和的乘积，而约束条件本身又要求下一时刻的状态满足雅可比矩阵与状态向量之积为零，由此便可以联立求解出用于满足约束条件所需的拉格朗日乘子，这一方程组的维度等于约束条件的数目。以一个受重力作用、约束在 $y=2x$ 直线上运动的质点为例，可以具体求解出拉格朗日乘子及对应的约束力，进而得到质点沿约束方向运动的新速度。

在实际系统中，约束条件的推导需要结合欧拉定理：已知某刚体的角速度以及其上某一固连点 $A$ 的速度，便可以推算出该刚体上任意其他点的速度。以两个刚体在某点通过关节连接、要求二者在该点平动速度保持一致为例，可以将约束条件写成关于两个刚体各自平动与转动速度的线性组合，并进一步整理出各自对应的雅可比矩阵，其中还会用到把叉乘运算表示为矩阵乘法的"叉乘矩阵"。当然，实际实现远比这一基本原理复杂，还需要考虑求解方程组的具体算法、约束条件的数值稳定化处理，以及用于表征刚体之间接触关系的单边约束条件等诸多方面。

以一台由四个刚体和三个转动关节构成的三轴机器人为例，其状态向量维度为二十四，质量矩阵维度为二十四乘二十四，雅可比矩阵维度为十五乘二十四，拉格朗日乘子向量维度为十五，由此形成一个维度为十五、用于求解拉格朗日乘子的方程组，而后续的正向动力学计算则可以独立于此、无需额外的矩阵求逆即可完成。规模更大的系统会带来更高的方程组维度，例如某林业机械模型包含五十个刚体和五十七个约束条件，对应的状态向量维度为三百，拉格朗日乘子向量维度为二百八十五，雅可比矩阵和质量矩阵的维度也相应增大；再如一辆车辆的动力学模型，其约束条件涉及车身与空气动力学、车轮悬挂、车轮与路面接触、制动系统、传动系统、转向系统、驾驶员等诸多子系统，涵盖诸如基于牵引杆模型与PID控制器的转向、按车轮分配的制动力矩、发动机特性曲线与差速器及变速箱、转向传动与阿克曼转向几何、经验轮胎模型等具体环节。综合起来，刚体动力学仿真的一次典型循环包括四个环节：由状态向量出发，依据约束条件方程求解约束力（即所谓的逆动力学环节，由运动推算约束力），再依据全部力（外力与约束力）推算下一时刻的运动状态（即正向动力学环节），如此循环推进，构成整个仿真过程在时间上的演化。

\subsection{跨领域仿真方法}

生产的网络化对仿真本身也提出了新的要求。在传统的生产环境中，机床、工件、工人、传感器与通信网络分属不同的专业领域，也各自被单独建模和仿真：机械专业计算力学行为，电气专业计算控制回路，物流专业计算节拍与产能，彼此之间几乎没有交集。而在互联的生产系统中，各要素之间的\textbf{相互作用}恰恰成为系统整体行为的决定因素，孤立的单领域仿真已不足以刻画真实表现，因此\textbf{必须进行整体性的仿真}——把此前彼此分离的机器、工人、传感器、通信等诸多过程\textbf{汇聚到同一个仿真环境之中}统一考察。这也正是以下几类跨领域仿真方法出现的动因。

具体而言，从系统仿真的视角出发，一个生产系统中以下几个方面都是相关的、需要被一并纳入模型：
\begin{itemize}
  \item \textbf{运动学与动力学}：设备自身的运动能力与受力行为；
  \item \textbf{传感器与执行器}：系统感知环境、并作用于环境的接口；
  \item \textbf{通信基础设施}：各组件之间交换信息的通道及其带宽与延迟；
  \item \textbf{数据处理}：对所采集数据进行的计算、评估与决策；
  \item \textbf{控制与管理系统（Leitsysteme）}：单元层以上的调度与指令下达；
  \item \textbf{用户界面}：人向系统下达指令、系统向人反馈状态的方式；
  \item \textbf{工人}：作业人员的行为、节拍与判断；
  \item \textbf{环境}：工件、周边设施以及外部条件；
  \item \textbf{相邻单元与更上层的控制技术}：处在本单元之外、却与之持续交互的部分。
\end{itemize}
这些方面此前分属不同专业、各自建模，而互联的生产恰恰要求把它们置于同一个模型中共同考察。

除了针对具体物理领域的专门仿真方法之外，还存在若干种通用性更强、可以跨越不同物理领域使用的建模与仿真范式。

基于方程的仿真，是通过代数微分方程组（即由常微分方程与代数方程共同构成的系统）来描述系统与过程，其一般形式可以写成状态方程与输出方程两组关系式；求解这类方程组通常需要借助适当的（有时是迭代式的）数值方法，这类仿真方式在机电一体化技术系统的仿真领域应用广泛，典型的建模载体是脚本语言或方框图，MATLAB即为其代表之一。以借助MATLAB对一台UR10机器人进行位置控制为例，仿真循环中依次完成推进仿真时间步、获取当前关节位置与速度、由信号发生器给出目标设定点、通过控制器计算力矩、将力矩写回仿真器、并记录与可视化数据这一系列步骤。

信号导向的仿真，是通过\textbf{作用图（Wirkungsplan）}来描述系统与过程；这里所说的作用图，是指借助方块、加法节点与分支节点、并通过作用线连接起来，对系统内部作用流程的一种符号化表示。以方框图形式呈现的模型，等价于给出了该模型背后的方程组，\textbf{模型的输入输出量（即信号流向）必须被明确规定}，这种方式在控制工程领域尤为常见，模型的各个部分还可以被封装为独立的模块，从而使信号流向更加清晰、整体结构更加直观，Simulink即是这类工具的代表；同样以UR10机器人的位置控制为例，可以借助信号发生器、位置误差与速度反馈构成的比例-微分控制模块、以及代表机器人动力学的仿真模块，将各部分以方框图的形式连接起来，构成完整的闭环仿真。

面向对象的仿真，也称为物理化面向对象建模或网络建模，是通过若干相互连接（即节点）的元件来描述系统，物理上的关联通过势量与流量这两类物理量的表达式来刻画；与信号导向仿真不同的是，模型方程中的输入输出量并非预先规定，而是\textbf{随着各元件之间的连接关系自然浮现出来，连接本身传递的是物理量，不具有方向性，也无需事先对模型方程做数学层面的整理}，这一整理工作由仿真工具自动完成，Modelica即是这类描述语言的代表。以一台直流电机的跨领域Modelica模型为例，其模型由电阻、电感、直流电压源、接地、机电转换元件以及负载惯量共同构成，模型代码中通过一系列connect语句描述各元件之间的连接关系，整个模型天然地融合了电气回路与机电、机械部分。

事件离散的仿真，只对系统中依次发生的各类事件本身进行模拟，而不关心事件之间间隔的具体时间流逝过程：在离散过程的运行中出现的所有事件都会被重演，每一个事件都会触发一段依赖于该事件类型的事件处理程序；由此产生的输出，例如状态及状态转换的记录，还可以进一步用于统计诸如停留时长等其他特征量。这种方式还可以进一步区分时间与数值上连续或离散的不同组合。以一台UR10机器人执行抓取动作为例，从流程启动开始，依次经过机器人移动到抓取点、闭合夹爪、离开抓取点这几个环节，每个环节对应一个离散事件（机器人已移动、夹爪已闭合等），而这些事件之间机器人具体如何运动的连续过程，在这一层面的仿真中并不被关心。

在系统层面，基于模型的系统工程是系统工程的一个分支，其核心在于形式化地运用建模手段，服务于系统需求采集、设计、分析、验证与确认等一系列活动，这些活动始于概念设计阶段，并贯穿整个开发过程乃至后续的生命周期各阶段；系统工程本身则更关注确保各个组成部分协同配合、共同实现整体目标，是一种跨学科的方法，聚焦于开发早期阶段系统需求的界定与文档化、系统架构的构建，以及在兼顾运行、时间、测试、建造、成本与规划、培训与支持、报废处置等整体问题的前提下，对系统是否满足既定需求加以核查，SysML是这一领域应用最为广泛的建模语言之一。

\subsection{面向工业4.0的概念与应用场景}

在具体的工业应用层面，\textbf{数字化工厂是对一整套由数字模型、方法与工具构成的综合网络的统称}，这些要素通过一体化的数据管理相互整合，其目标是对现实工厂中与产品相关的各类结构、流程与资源进行\textbf{整体性的规划、评估与持续改进}；数字化工厂的核心目标在于加速并改进生产规划工作，并以"并行工程"的方式，使其与产品开发在时间上高度重叠，其应用范围覆盖从市场需求、设计、物料清单、项目管理、原型数字样机，到内外部物流、装配与制造工艺规划、生产设施规划、生产设施的装配与调试、投产管理、批量生产、销售与订单处理、外部物流、服务与维护等生产全流程的各个环节。

与之密切相关的是"虚拟调试"：狭义上，虚拟调试是指借助设施的仿真模型，对自动化系统在真正投入现场调试之前所进行的整体性预先测试；广义上，虚拟调试还包括在开发过程中，针对自动化系统的各个组件与部分功能，采用与具体任务相适配的仿真方法和仿真模型所开展的伴随性测试。虚拟调试的典型目标包括发现控制程序中的错误（例如逻辑错误、录入或复制错误、错误的使能与联锁信号）、测试操作界面、测试并验证设施的动态行为（例如满载与空载运行、调试模式、手动模式或自动模式）、考虑各类特殊情形（例如部件的进出料处理）、验证既定的设施参数（例如节拍时间、产量、可用性），以及有针对性地仿真各类故障情形（例如线缆断裂、传感器失效）；在硬件在环这一配置下，真实的控制器与真实的通信协议（例如Profibus、Profinet、EtherCAT）连接到运行在仿真计算机上的、包含传感器与执行器在内的设施仿真模型，从而在不接触真实设施的情况下完成上述各项测试。

在具体案例层面，可重构装配单元的规划与运行是仿真与数字孪生理念的一个典型应用：从期望产品出发，经由CAD模型及附加信息的自动化预处理，进入仿真与编程环节——依次生成装配序列、通过可视化编程转化为"技能"序列、进行单元仿真、并结合单元重构完成对候选方案的评估与工艺控制，最终交由具备反馈控制、力控制与视觉感知能力的实际装配单元执行，执行结果又反过来用于修正和迭代前述仿真与编程环节，形成一个闭环。与此类似的还有基于仿真的机器人工作单元优化，例如通过仿真不同的相机安装角度与仿真时长组合，以画面中可见的目标像素数量作为评价指标，寻找视野最佳的相机布置方案；还有面向手工工位的自动化改造项目，其中需要在多种候选机器人系统之间进行选择，选择过程依托于对人工作业流程的建模、目标状态的生成、规划、仿真与评估这一整套"行动规划"框架，并最终依据仿真评估结果确定合适的机器人系统。

除此之外，辅助系统也是数字孪生理念的重要应用场景之一，例如借助增强现实技术，将设备的数字模型叠加到真实的车床画面之上，辅助操作与维护人员的日常工作。而在自动驾驶领域，数字孪生同样扮演着不可或缺的角色：车辆需要依托对周边环境（道路、树木、其他车辆等）的精确三维重建与特征点匹配来实现感知与定位，这类感知与仿真能力正是自动驾驶系统得以实现的基础之一。